\chapter{绪论}\label{chap:introduction}

\section{研究背景与意义}

同步定位与地图构建（Simultaneous Localization and Mapping, SLAM）是移动机器人、无人机以及增强现实/虚拟现实设备实现自主导航与环境感知的关键技术\cite{thrun2005probabilistic}。在未知或弱先验环境中，系统需依靠机载传感器持续估计自身位姿并构建环境地图，从而为路径规划、任务执行和人机交互提供基础支撑。Cadena 等人\cite{cadena2016past}在其里程碑式的综述中系统回顾了 SLAM 研究三十余年的发展历程，将其划分为经典时代（1986--2004）、算法改进时代（2004--2015）与鲁棒感知时代（2015 至今），指出 SLAM 已从单一的概率推断问题演进为融合几何、语义和深度学习的综合性感知框架。随着视觉传感器的成本持续下降、成像质量提升以及嵌入式计算能力增强，视觉 SLAM 逐渐成为主流方案\cite{gao2017vslam14}，并在空间探测、仓储物流、家庭服务与移动终端等场景中得到广泛应用。

从应用场景看，SLAM 已成为多类智能系统的基础能力。在自动驾驶与配送机器人中，SLAM 为车辆提供高精度定位与动态环境建图，支撑路径规划和安全决策；在无人机与室内巡检中，SLAM 使平台在 GPS 不可用的环境中实现稳定导航\cite{forster2014svo}；在 AR/VR 设备与移动终端中，SLAM 支撑实时空间理解与虚实融合，直接决定交互体验\cite{klein2007parallel}；在仓储、安防与服务机器人场景，SLAM 则提供可重复构建的语义地图与任务级定位能力。图~\ref{fig:slam_applications} 展示了 SLAM 技术在不同领域的典型应用场景。

\begin{figure}[htbp]
    \centering
    \begin{subfigure}[b]{0.48\textwidth}
        \centering
        \includegraphics[width=\textwidth]{Img/scene1_auto_drive_car.png}
        \caption{自动驾驶汽车}
        \label{fig:slam_app_auto_drive}
    \end{subfigure}
    \hfill
    \begin{subfigure}[b]{0.48\textwidth}
        \centering
        \includegraphics[width=\textwidth]{Img/scene2_delivery_robot.png}
        \caption{配送机器人}
        \label{fig:slam_app_delivery}
    \end{subfigure}

    \vspace{0.5em}

    \begin{subfigure}[b]{0.48\textwidth}
        \centering
        \includegraphics[width=\textwidth]{Img/scene3_drone.png}
        \caption{无人机室内导航}
        \label{fig:slam_app_drone}
    \end{subfigure}
    \hfill
    \begin{subfigure}[b]{0.48\textwidth}
        \centering
        \includegraphics[width=\textwidth]{Img/scene4_AR_VR.png}
        \caption{增强/虚拟现实设备}
        \label{fig:slam_app_ar_vr}
    \end{subfigure}

    \vspace{0.5em}

    \begin{subfigure}[b]{0.48\textwidth}
        \centering
        \includegraphics[width=\textwidth]{Img/scene5_industrial_mobile_robot.png}
        \caption{仓储搬运机器人}
        \label{fig:slam_app_warehouse}
    \end{subfigure}
    \hfill
    \begin{subfigure}[b]{0.48\textwidth}
        \centering
        \includegraphics[width=\textwidth]{Img/scene6_security_patrol_robot.png}
        \caption{安防巡检机器人}
        \label{fig:slam_app_security}
    \end{subfigure}

    \vspace{0.5em}

    \begin{subfigure}[b]{0.48\textwidth}
        \centering
        \includegraphics[width=\textwidth]{Img/scene7_humanoid_service_robot.png}
        \caption{家庭服务机器人}
        \label{fig:slam_app_service}
    \end{subfigure}
    \hfill
    \begin{subfigure}[b]{0.48\textwidth}
        \centering
        \includegraphics[width=\textwidth]{Img/scene8_lunar_rover.png}
        \caption{空间探测月球车}
        \label{fig:slam_app_space}
    \end{subfigure}

    \bicaption{SLAM 技术的典型应用场景}{Typical application scenarios of SLAM technology}
    \label{fig:slam_applications}
\end{figure}

SLAM 的发展历史可以追溯到 1980 年代末和 1990 年代初的概率机器人学研究。Thrun 等人\cite{thrun2005probabilistic}在其专著《Probabilistic Robotics》中系统阐述了将概率推断方法应用于机器人感知与决策的理论框架，为 SLAM 的数学建模奠定了基础。这一时期的核心贡献在于将地图构建与定位问题统一在一个概率框架下，特别是 Smith 和 Cheeseman~\cite{smith1988estimating} 提出的状态估计理论，以及 Durrant-Whyte~\cite{durrant1987uncertain} 对不确定性建模的深入探讨。随着 Smith 等人~\cite{smith1990estimating} 明确描述了地标位置与机器人位姿之间的统计相关性，SLAM 作为一个独立的科学问题正式确立。Durrant-Whyte 和 Bailey~\cite{durrantwhyte2006slam} 在其经典的双篇综述教程中（Part I 与 Part II\cite{bailey2006simultaneous}），分别对 SLAM 的概率基础与计算方法进行了系统梳理：Part I 阐述了基于 EKF、粒子滤波与信息滤波的递推贝叶斯框架，深入分析了 SLAM 问题的收敛特性与一致性条件；Part II 则讨论了数据关联、闭环检测和多机器人协作等关键工程问题。此后基于扩展卡尔曼滤波（EKF-SLAM）的方法成为早期最为主流的实现方案。这一阶段的关键认识是地标误差与位姿误差存在强相关性，地图无法拆分为独立估计问题，因此必须在联合状态中统一建模。与此同时，EKF-SLAM 的协方差矩阵随地标数量 $N$ 呈 $O(N^2)$ 增长，计算与存储开销随规模迅速上升，直接推动了后续的 FastSLAM 与图优化框架的发展。进入 2000 年代，随着图优化理论与稀疏线性代数领域的成熟，Kaess 等人先后提出了增量式平滑与建图方法 iSAM\cite{kaess2008isam} 和 iSAM2\cite{kaess2011isam2}，开创性地利用贝叶斯树数据结构实现了因子图的增量更新，使得大规模 SLAM 问题的在线求解成为可能。K{\"u}mmerle 等人\cite{kummerle2011g}提出的通用图优化框架 g$^2$o 则为后续众多 SLAM 系统提供了高效的后端优化引擎，支持多种求解器与鲁棒核函数。2010 年代以来，视觉 SLAM 迅速普及，单目、双目与视觉惯性系统相继成熟，ORB-SLAM\cite{mur2015orb}、LSD-SLAM\cite{engel2014lsd}、SVO\cite{forster2014svo}、OKVIS\cite{leutenegger2015keyframe}、VINS-Mono\cite{qin2018vins} 与 ROVIO\cite{bloesch2015robust} 等经典框架推动了系统在实时性、鲁棒性与多传感器融合方面的进步，ORB-SLAM2/3 进一步扩展到双目、RGB-D 与视觉惯性多地图场景\cite{mur2017orb,campos2021orb}。近年来，深度学习与语义感知开始与 SLAM 融合，使得系统在动态场景与弱纹理环境中的表现进一步增强，同时也对后端优化的算力提出了更高要求。

典型视觉 SLAM 体系通常包含传感器数据获取、前端跟踪、后端优化、回环检测与建图更新等模块，其流程如图~\ref{fig:visual_slam_pipeline} 所示。前端负责特征提取与匹配并给出初始位姿估计，后端通过非线性优化进一步消除累积误差，回环检测与建图模块则保证全局一致性与地图质量。由于视觉测量噪声、光照变化与动态遮挡等因素的存在，前端估计易产生漂移，因此位姿优化成为决定定位精度与系统稳定性的核心环节。

\begin{figure}[htbp]
    \centering
    \resizebox{0.92\linewidth}{!}{\begin{tikzpicture}[node distance=14mm,>=Stealth,rounded corners,thick]
    \tikzstyle{block}=[draw,minimum width=22mm,minimum height=8mm,align=center]
    \tikzstyle{data}=[draw,minimum width=18mm,minimum height=8mm,align=center]

    \node[block] (sensor) {传感器\\数据获取};
    \node[block, right=of sensor] (frontend) {前端\\跟踪与匹配};
    \node[block, right=of frontend] (backend) {后端\\位姿优化};
    \node[block, right=of backend] (loop) {回环\\检测};
    \node[block, below=of backend] (mapping) {建图\\与更新};

    \draw[->] (sensor) -- (frontend);
    \draw[->] (frontend) -- (backend);
    \draw[->] (backend) -- (mapping);
    \draw[->] (loop) -- (backend);
    \draw[->] (backend) -- (loop);
    \draw[->] (mapping) -| (frontend);
\end{tikzpicture}
}
    \bicaption{经典视觉 SLAM 系统流程}{Classic visual SLAM pipeline}
    \label{fig:visual_slam_pipeline}
\end{figure}

在视觉 SLAM 的演进历程中，后端状态估计主要形成了以扩展卡尔曼滤波（EKF）为代表的滤波法和以束调整（Bundle Adjustment, BA）为代表的非线性优化法两大路径。早期的后端方案多采用基于马尔科夫假设的滤波框架，如经典 EKF-SLAM\cite{smith1990estimating}，其核心在于通过预测与更新步骤实时维护状态向量及其协方差矩阵；随后，Mourikis 等人提出的多状态约束卡尔曼滤波（MSCKF）\cite{mourikis2007multi}进一步提升了视觉惯性系统的实时性。滤波法的优势在于计算资源消耗相对可控，通过边际化历史状态实现了极高的处理频率；然而，其固有的线性化误差往往导致系统估计的不一致性，且在处理大规模闭环或长距离漂移时，缺乏重新线性化的机制，容易陷入局部最优。相比之下，基于图优化的非线性优化法\cite{kummerle2011g}逐渐成为当前主流趋势，该方法将定位与建图转化为超定方程的最小二乘问题，通过高斯—牛顿或列文伯格—马夸尔特等算法迭代求解\cite{triggs1999bundle}。优化法的核心优势在于允许对历史状态进行重新线性化，配合滑动窗口或关键帧机制\cite{leutenegger2015keyframe}，能显著提升全局定位精度；然而，这种精度增益是以巨大的计算开销为代价的，优化过程涉及复杂的雅可比矩阵计算、大规模海森矩阵累积以及稀疏线性方程组求解，尽管稀疏性可利用舒尔补进行加速，但随着路标点与关键帧规模的迅速扩张，计算复杂度仍呈指数级增长，形成显著的后端性能瓶颈。在高频视觉输入与实时性需求下，这种计算密集的特性使得后端优化在移动端和嵌入式平台上难以满足低延时要求。传统的通用处理器（CPU）或图形处理器（GPU）虽然具备极高的编程灵活性，但受限于功耗与能效比，难以在小型无人机、可穿戴 AR 设备等对功耗极度敏感的场景中长期稳定运行。因此，针对后端优化算法中矩阵运算的底层特征进行硬件架构级的加速设计，已成为突破 SLAM 系统实时性瓶颈的关键方向。

与通用处理器相比，现场可编程门阵列（Field-Programmable Gate Array, FPGA）具备高度并行、确定性时延与可定制存储层级的优势，能够以更低能耗实现稀疏矩阵运算、雅可比流水线与舒尔补求解等核心算子，从而为 SLAM 后端优化提供更高的实时性保障\cite{huangkun2025vslam}。通过软硬协同的设计方式，将计算密集的位姿优化环节迁移至 FPGA，同时在处理器端保留系统控制与轻量更新逻辑，可在保持精度的同时显著降低延时与功耗。这不仅有助于推动视觉 SLAM 在资源受限平台上的落地，也为高可靠、自主运行的智能系统提供关键支撑。

\section{国内外研究现状}

针对视觉 SLAM 及其后端优化的加速研究，国内外学者从算法改进与硬件架构设计两个维度开展了大量工作。本节将从软件算法演进、国外研究现状以及国内研究进展三个方面进行综述。

\subsection{系统框架与软件算法演进}

视觉 SLAM 算法的演进是一个多种理论不断交叉融合的过程，从早期的滤波机制到现代的非线性优化，再到如今引入深度学习模型，其系统框架在精度、鲁棒性以及计算复杂度上不断寻求平衡。本节将从前端数据关联方式、后端状态估计框架以及深度学习的融合三个维度，对视觉 SLAM 的软件算法演进进行阐述。

\subsubsection{视觉前端数据关联}

视觉前端的主要任务是从图像流中提取特征并建立帧间数据关联，进而为后端提供初始的相机运动估计。根据数据关联方式的不同，前端主要分为间接法、直接法和半直接法三大类。

间接法（特征点法）是目前最为成熟和广泛应用的前端方案。该方法首先从图像中提取具有代表性的关键点并计算相应的描述子，随后通过描述子匹配建立相邻帧之间的 2D-2D 或 2D-3D 数据关联，最后通过最小化重投影误差（Reprojection Error）求解相机位姿。在特征提取器的发展历程中，Lowe\cite{lowe2004distinctive} 于 2004 年提出的尺度不变特征变换（SIFT）是里程碑式的工作，该算法通过构建高斯差分尺度空间检测关键点，并以关键点邻域梯度直方图作为 128 维描述子，具有良好的旋转、尺度与仿射不变性，但其计算开销较大（单帧需数百毫秒），难以满足实时性要求。Bay 等人\cite{bay2008surf} 提出的加速稳健特征（SURF）利用积分图像与 Haar 小波近似替代 SIFT 中的高斯卷积，将计算速度提升约 3 倍，但仍无法满足嵌入式平台的帧率需求。为此，Rublee 等人\cite{rublee2011orb} 在 2011 年提出了 ORB（Oriented FAST and Rotated BRIEF）描述子，结合 Rosten 和 Drummond\cite{rosten2006machine} 设计的高速 FAST 角点检测器与旋转 BRIEF 二进制描述子，在保持良好匹配性能的同时将单帧计算压缩到毫秒级别，成为当前实时 SLAM 系统中应用最为广泛的手工特征。间接法的优势在于对光照变化和图像模糊具有较强的鲁棒性，且提取的几何特征便于后端进行优化与回环检测。

在系统层面，Klein 和 Murray\cite{klein2007parallel} 于 2007 年提出的 PTAM（Parallel Tracking and Mapping）开创性地将跟踪与建图分离到两个并行线程，首次在增强现实场景中实现了实时的关键帧式 BA 优化，标志着 SLAM 从滤波范式向优化范式的重要转折。此后，Mur-Artal 等人\cite{mur2015orb}提出的 ORB-SLAM 将 ORB 特征、共视图（Covisibility Graph）与 G{\'a}lvez-L{\'o}pez 和 Tard{\'o}s\cite{galvez2012bags} 提出的词袋（Bag of Binary Words, DBoW2）回环检测融为一体，构建了包含跟踪、局部建图与闭环三线程的完整系统架构。ORB-SLAM2\cite{mur2017orb} 将支持范围从单目扩展到双目与 RGB-D 相机，ORB-SLAM3\cite{campos2021orb} 进一步集成了视觉惯性紧耦合与多地图管理能力，成为当前功能最完备的开源特征点法 SLAM 系统。然而，特征提取和匹配过程极为耗时，往往构成了前端计算资源的瓶颈，尤其在高分辨率图像或高帧率场景下，ORB 提取可占据前端超过 60\% 的计算时间。

直接法（Direct Method）的提出旨在跳过耗时的特征提取环节。它基于灰度不变性假设（即空间中同一点在不同视角下的像素灰度保持不变），通过直接最小化光度误差（Photometric Error）来估计相机运动。直接法能够利用图像中的像素点（包括边缘和无明显角点区域），在弱纹理环境中表现出更好的鲁棒性。Newcombe 等人\cite{newcombe2011dtam}提出的 DTAM（Dense Tracking and Mapping）是最早的实时稠密直接法系统之一，通过最小化全图光度误差同时恢复相机运动与稠密深度图，但依赖 GPU 并行计算。Engel 等人\cite{engel2014lsd}提出的 LSD-SLAM（Large-Scale Direct Monocular SLAM）将直接法与半稠密重建结合，在单目 CPU 上实现了大规模环境的实时建图。此后，Engel 等人\cite{engel2017direct}进一步提出了 DSO（Direct Sparse Odometry），放弃了 LSD-SLAM 中的稠密重建策略，改为在图像中选取具有强梯度的稀疏像素点进行联合光度优化，同时显式建模几何参数（逆深度）与光度参数（曝光时间与暗角），在保持直接法灵活性的同时显著提升了精度与效率。尽管直接法省去了特征计算，但由于局部光度梯度的非凸性，其对相机的剧烈运动和光照突变极为敏感，容易陷入局部最优。

为了结合上述两者的优势，半直接法（Semi-direct Method）应运而生。半直接法通常仅提取 FAST 角点但不计算描述子，利用光流法或直接法追踪这些特征点，随后再结合特征对齐提升定位精度。它在一定程度上避免了直接法的非凸性问题，同时免去了特征点法冗长的描述子匹配开销。Forster 等人\cite{forster2014svo}于 2014 年提出的 SVO（Semi-direct Visual Odometry）是该领域的标杆算法，其将稀疏图像对齐（Sparse Image Alignment）与特征对齐（Feature Alignment）级联，凭借极高的前端运行帧率（在标准测试中可达 300~fps 以上）和良好的精度，成为微型飞行器（MAV）等计算受限平台的理想框架。SVO 2.0\cite{forster2017svo} 进一步扩展了多相机支持与边缘特征利用，显著提升了系统在宽基线与大视场角场景下的鲁棒性。

\subsubsection{后端状态估计框架}

视觉 SLAM 的后端负责消除前端的累积漂移，通过融合多源测量（如 IMU、里程计等）维持全局一致的地图。按照状态估计类型的不同，主要分为滤波法和优化法两大路径。

滤波法（Filter-based）主要基于马尔科夫假设，利用贝叶斯滤波框架在当前时刻更新系统状态与协方差矩阵。早期的单目 MonoSLAM\cite{davison2007monoslam} 采用了扩展卡尔曼滤波（EKF）框架，将相机位姿与路标点共同纳入状态向量进行更新。Davison 等人的贡献在于首次证明了单目相机在未知环境中能够实时完成概率性的自主建图，但由于状态向量和协方差矩阵的大小随路标点的增加而膨胀，其计算量呈平方级增长（$O(N^2)$），严重限制了地图的大规模扩展。为解决此问题，Mourikis 和 Roumeliotis\cite{mourikis2007multi} 提出了 MSCKF（Multi-State Constraint Kalman Filter），其核心思想是在状态向量中仅维护滑动窗口内的历史位姿，将针对同一路标的多帧观测投影到零空间消除路标状态，使计算复杂度仅与窗口大小相关而非路标数量。MSCKF 极大地提升了系统的实时性，成为了视觉惯性里程计（VIO）领域的经典方案，并广泛应用于 Google ARCore 等工业级产品。Geneva 等人\cite{geneva2020openvins}在 ICRA 2020 上发布的 OpenVINS 开源平台系统整合了 MSCKF 的多种变体与扩展（包括 SLAM 特征、FEJ 一致性保证、在线标定等），为 VIO 算法的研究与对比评估提供了统一基准。Bloesch 等人\cite{bloesch2015robust}提出的 ROVIO（Robust Visual Inertial Odometry）则将直接法光度误差引入 EKF 更新步骤，避免了特征提取与匹配的开销，在计算资源极为有限的场景下展现出良好的实时性。

随着图优化理论与稀疏矩阵线性代数求解器领域的成熟，优化法（Optimization-based）逐渐取代滤波法成为主流框架。Kaess 等人\cite{kaess2008isam}于 2008 年提出的 iSAM（Incremental Smoothing and Mapping）首次实现了因子图的增量式 QR 分解更新，避免了每次新增观测时对整个信息矩阵的重新分解。此后，iSAM2\cite{kaess2011isam2} 引入贝叶斯树（Bayes Tree）数据结构，实现了增量式变量重排序与流体式重新线性化（Fluid Relinearization），使得仅受新观测影响的局部子图被更新，极大提升了大规模地图的在线求解效率。K{\"u}mmerle 等人\cite{kummerle2011g}提出的通用图优化框架 g$^2$o（General Graph Optimization）为后续众多 SLAM 系统提供了高效的后端优化引擎，其采用稀疏 Cholesky 分解求解正规方程，支持多种鲁棒核函数与求解策略，已被 ORB-SLAM、LSD-SLAM 等系统广泛采用。优化法将所有或一段滑动窗口内的历史轨迹及路标点构建为一个超大的非线性最小二乘图，通过束调整（Bundle Adjustment, BA）全局解算位姿与路标。Triggs 等人\cite{triggs1999bundle}在其经典的综述中系统总结了 BA 的理论基础与实现方法，指出 BA 本质上是一个在流形上的稀疏非线性最小二乘问题，其海森矩阵（Hessian Matrix）天然具有箭头形稀疏结构，可通过舒尔补（Schur Complement）消元将路标变量边际化，从而将求解规模从 $(N_p + N_m)$ 降至 $N_p$（$N_p$ 为位姿数，$N_m$ 为路标数），这一思想构成了本文硬件加速器设计的算法基础。

在视觉惯性融合方面，IMU 预积分理论的提出是重要的里程碑。Lupton 和 Sukkarieh\cite{lupton2012visual} 于 2012 年首次提出了 IMU 预积分（Preintegration）的概念，将相邻关键帧之间的高频 IMU 测量预先积分为相对运动约束，避免了偏置更新时对所有 IMU 数据的重新积分。Forster 等人\cite{forster2017onmanifold}在此基础上进一步发展了流形上的预积分理论（On-Manifold Preintegration），通过在 $SO(3)$ 流形上直接进行旋转积分，消除了欧拉角或四元数参数化带来的奇异性问题，并推导了协方差传播与偏置修正的解析公式，成为当前 VIO 系统（如 VINS-Mono、OKVIS、ORB-SLAM3）普遍采用的标准方法。

当前主流系统如 ORB-SLAM 系列、OKVIS\cite{leutenegger2015keyframe} 以及 VINS-Mono\cite{qin2018vins} 均采用优化法。其中，Leutenegger 等人\cite{leutenegger2015keyframe}提出的 OKVIS 是最早的基于关键帧的视觉惯性非线性优化系统之一，通过在滑动窗口中联合优化相机位姿、速度、IMU 偏置与路标点，建立了紧耦合 VIO 的基本范式。Qin 等人\cite{qin2018vins}提出的 VINS-Mono 进一步将滑动窗口机制与预积分技术相结合，创新性地设计了边缘化（Marginalization）策略将早期状态信息转化为先验融入系统，并引入了在线时序标定\cite{qin2018online}以处理视觉与惯性传感器之间的时间偏移，在 EuRoC\cite{Burri25012016} 等权威数据集上展现了优异的精度表现。此外，Fan 等人\cite{fan2024schurvins}在 CVPR 2024 上提出的 SchurVINS 通过将舒尔补消元与误差状态卡尔曼滤波（EKF）增量更新相结合，以极低的计算代价实现了接近全 BA 优化的精度，为嵌入式平台上的 VIO 部署提供了轻量化解决方案，本文的硬件加速架构正是基于该算法框架展开设计的。

\subsubsection{深度学习与视觉 SLAM 的融合}

近年来，深度学习技术正在逐步渗透到 SLAM 的各个模块中，从特征提取与匹配到深度估计与端到端位姿回归，深层神经网络通过数据驱动的方式展现出了应对复杂动态场景、弱纹理区域以及剧烈光照变化的强大能力。

在前端特征提取层面，DeTone 等人\cite{detone2018superpoint}于 2018 年提出的 SuperPoint 网络采用自监督训练策略，通过单应性变换（Homographic Adaptation）从合成几何图形迁移学习到真实图像，能够在单次前向推理中同时输出关键点位置及其 256 维描述子。相比传统的 FAST/ORB 提取器，SuperPoint 在光照变化、视角变换与运动模糊条件下展现出显著更强的重复性与匹配稳定性。在此基础上，Sarlin 等人\cite{sarlin2020superglue}于 CVPR 2020 提出的 SuperGlue 利用图神经网络（Graph Neural Network）与注意力机制对关键点集合进行联合推理，通过可微分的 Sinkhorn 最优传输算法建立跨图像的特征对应关系，显著提升了宽基线、重复纹理和弱纹理场景下的匹配精度。Sun 等人\cite{sun2021loftr}在 CVPR 2021 进一步提出了 LoFTR（Local Feature Matching with Transformers），该方法完全摒弃了关键点检测步骤，利用 Transformer 的全局感受野在粗粒度特征图上建立稠密对应，再通过精细化模块提取亚像素级匹配。LoFTR 在室内弱纹理环境（如白墙、走廊）中表现尤为突出，为无纹理区域的视觉定位提供了新的思路。

在深层多视图几何理解方面，Teed 和 Deng\cite{teed2020raft} 提出了基于全对偶场循环架构的 RAFT（Recurrent All-Pairs Field Transforms）用于准确预估光流。RAFT 构建了 4D 全对偶相关体（Correlation Volume），并通过 GRU 循环单元迭代更新光流场，在 Sintel 和 KITTI 基准上取得了突破性的精度提升。到了系统层面，Teed 和 Deng\cite{zhu2020droid} 进一步提出了端到端的 DROID-SLAM，利用循环 GRU 模块迭代更新相机位姿与像素级对应场，在前端跟踪阶段采用稠密 BA 层（Dense BA Layer）通过可微分的 Gauss-Newton 求解实时优化位姿与深度，后端则借助全局 BA 消除累积漂移。DROID-SLAM 在 TartanAir、EuRoC 与 TUM-RGBD 等多种基准上均达到了与经典几何方法可比甚至超越的精度，特别是在大运动与严重模糊条件下展现出了显著优于传统方法的鲁棒性。

尽管深度学习 SLAM 展现出了高鲁棒性与语义感知潜力，但其依赖高性能 GPU 运算的特性（DROID-SLAM 在 NVIDIA RTX 3090 上仍仅能以约 10~fps 运行），极大地限制了其在功耗和算力受限的边缘芯片平台上的大规模部署。此外，深度学习方法对训练数据分布的依赖使其泛化能力仍有待提升。因此，在当前的嵌入式 SLAM 系统中，基于传统几何方法的优化后端仍然是保证实时性与可靠性的主流选择。

\subsubsection{后端优化的计算瓶颈分析}

综上所述，虽然基于优化法（BA）的 SLAM 在精度与鲁棒性方面表现优异，但在嵌入式系统中它却带来了沉重的计算负担。莫霄睿等人\cite{mo2024ba_review}在其综述中系统分析了 BA 优化芯片的研究现状，指出一次典型的后端 BA 优化需要经历雅可比矩阵求导（Jacobian Computation）、海森矩阵累积（Hessian Accumulation）、舒尔补消元（Schur Complement Elimination）以及线性方程组求解（Linear System Solving）等多个计算密集阶段。以滑动窗口大小为 4、路标点数量为 200 的典型 VIO 场景为例，单次优化迭代需要完成约 $4 \times 200 \times 2 = 1600$ 次雅可比计算（考虑双目），累积 $6 \times 6$ 位姿块与 $3 \times 3$ 路标块共计数千个海森子矩阵，并最终求解 24 维的舒尔补正规方程组。Gan 等人\cite{gan2021eudoxus}在 HPCA 2021 上发表的 Eudoxus 工作中，通过对自动驾驶定位系统的系统性性能剖析（Profiling），量化揭示了后端优化在总计算时间中占比高达 70\%--90\% 的事实，其中舒尔补与矩阵求解是最主要的热点函数。这一复杂的矩阵迭代过程在移动终端（如基于 ARM 处理器的无人机或 AR 设备上）产生了几十到上百毫秒量级的延迟，成为制约系统实时性与能效比的瓶颈。

针对 SLAM 的高昂算力需求，学术界首先在前端探索了专用领域的硬件加速方案。Fang 等人\cite{fang2017fpga_orb}早在 2017 年就在 ICFPT 会议上展示了基于 FPGA 的 ORB 特征提取加速器，将 FAST 角点检测与 BRIEF 描述子计算整合到统一流水线中。此后，Liu 等人设计的 eSLAM\cite{liu2019eslam} 在 DAC 2019 上提出了面向完整 ORB-SLAM 前端的节能加速器，通过 FPGA 全流水线成功实现了高速特征检测、描述子计算与帧间匹配的端到端加速，相比 ARM Cortex-A9 实现了 3 倍加速与 61\% 的能耗节省。Fang 等人\cite{fang2019fpga}则在 IEEE TCAS-I 上发表了针对视觉 SLAM 前端的系统化加速架构，将特征提取、光流跟踪与位姿估计集成为统一的硬件流水线。Zhang 等人\cite{zhang2023fpga_feature}在 2023 年进一步提出了同时支持特征提取与追踪的 FPGA 加速器，通过并行化 Lucas-Kanade 光流计算实现了更高的前端吞吐率。然而，当前端的处理能力得到了成倍的提升后，原本就吃紧的后端优化阶段面临着更加密集的特征点输入约束。这意味着必须突破通用处理器的计算瓶颈，针对基于 BA 的稀疏矩阵运算法则设计更深层级与更大规模数据的专用处理架构，方能在低功耗平台满足 SLAM 复杂算法迭代的需要。因此，将目光聚焦于通过软硬协同或者底层算子加速来提升 SLAM 后端运算效率，已成为了突破嵌入式领域自动导航系统性能天花板的必由之路。

\subsection{国外研究现状}

国外顶尖高校与研究机构在 SLAM 硬件加速领域起步较早，研究重点已从早期的前端特征提取加速逐渐转向后端非线性优化与稀疏矩阵求解器的深度定制化设计。根据目标硬件平台的不同，相关工作可分为基于 FPGA 的前端视觉加速器、基于 FPGA 的后端优化加速器、基于 ASIC 的专用芯片流片以及基于 GPU 的并行优化等方向。

\subsubsection{基于 FPGA 的前端视觉加速器}

SLAM 前端的计算负载主要集中在特征提取与匹配、光流跟踪以及立体深度估计等环节。由于这些操作具有高度的数据级并行性与规则的访存模式，非常适合在 FPGA 上以流水线方式实现。

在特征提取加速方面，除前文提及的 Fang 等人\cite{fang2017fpga_orb}的 ORB 加速器和 Liu 等人\cite{liu2019eslam}的 eSLAM 系统外，清华大学电子工程系 Xu 等人\cite{xu2020cnn_fpga}于 FCCM 2020 上提出了基于 CNN 的 FPGA 特征点提取加速器，将轻量级卷积神经网络部署到嵌入式 FPGA 上，替代传统的手工特征（如 ORB/FAST），在保持与 SuperPoint 相当的检测精度的同时，实现了实时推理速度。该工作展示了深度学习特征在嵌入式 FPGA 上的可行性，为 SLAM 前端从手工特征向学习型特征的迁移提供了硬件基础。Han 和 Liu\cite{han2025orb_fpga}在 ICFPT 2025 上进一步提出了采用循环缓冲架构的 ORB 特征提取 FPGA 加速器，专门面向双目视觉 SLAM 场景，通过时间复用策略在单一硬件实例上交替处理左右目图像，有效降低了资源占用。

在角点检测与特征跟踪方面，Gu 等人\cite{gu2022harris}于 IROS 2022 上提出了面向视觉惯性里程计（VIO）的高能效 Harris 角点检测 FPGA 加速器。该工作来自清华大学精密仪器系，将 Harris 角点检测算法的梯度计算、自相关矩阵累积与非极大值抑制三个阶段整合到统一的像素级流水线中，在 Xilinx Zynq FPGA 上实现了实时 VIO 处理。在光流跟踪方面，Jang 等人\cite{jang2009klt}早在 2009 年就在 ITSC 会议上报告了金字塔 KLT（Kanade-Lucas-Tomasi）光流跟踪器的硬件实现，通过层次化的多分辨率光流计算与流水线化的迭代求解，实现了面向驾驶辅助系统的实时特征跟踪，为后续 SLAM 前端的光流硬件加速奠定了基础。Blachut 和 Kryjak\cite{blachut2022lk_fpga}在 2022 年进一步实现了多尺度 Lucas-Kanade 与 Horn-Schunck 光流算法在 FPGA 上的高效部署，通过像素级流水线与多分辨率金字塔并行计算，在 4K 分辨率视频流上实现了实时光流估计。Zhang 等人\cite{zhang2023fpga_feature}在 2023 年则进一步提出了同时支持特征提取与 Lucas-Kanade 光流追踪的 FPGA 加速器，通过并行化光流计算实现了更高的前端吞吐率。

在立体视觉与深度估计方面，双目立体匹配是 SLAM 系统获取稠密深度信息的重要途径。Banz 等人\cite{banz2010sgm}于 2010 年在 SAMOS 会议上首次系统地报告了半全局匹配（Semi-Global Matching, SGM）算法的 FPGA 实时实现方案，通过沿多个方向并行传播代价体（Cost Volume）并设计片上缓存策略来满足 SGM 算法的高带宽需求，实现了实时的稠密视差图计算。Shrivastava 等人\cite{shrivastava2020sgm}在 FPL 2020 上进一步提出了通过松弛 SGM 中路径间数据依赖关系来提升并行度的方法，在不显著降低视差图质量的前提下，将 FPGA 上的立体匹配吞吐率提升了 2 倍以上。

近年来，基于深度学习的光流估计模型（如 RAFT\cite{teed2020raft}）在精度上取得了突破性进展，其 FPGA 加速也成为新的研究热点。中山大学 Ling 等人\cite{ling2022flowacc}于 DATE 2022 上提出了 FlowAcc，这是首个面向基于 DNN 的光流估计的实时 FPGA 加速器，通过模型压缩与定制化数据流设计，在保持与 GPU 实现相当精度的同时实现了低功耗实时推理。同一团队的 Yan 等人\cite{yan2023optflow}在 Journal of Systems Architecture 2023 上发表了该工作的期刊扩展版本，进一步优化了加速器架构与量化策略。西安交通大学的 Jia 等人\cite{jia2025raft_fpga}于 ISCAS 2025 上专门针对 RAFT 的循环全对偶场变换（Recurrent All-Pairs Field Transforms）架构设计了 FPGA 加速器，通过特征相关体（Correlation Volume）的分块计算与 GRU 迭代单元的硬件映射，实现了端到端的实时光流估计。复旦大学的 Li 等人\cite{li2024optflow_slam}在 A-SSCC 2024 上则提出了面向 SLAM 场景的基于光流的 FPGA 加速器，创新性地利用帧间相似性（Inter-Frame Similarity）减少冗余计算，并采用相关性引导的混合精度流更新策略（Correlation-Guided Mixed-Precision Flow Update），在保持定位精度的同时显著降低了计算量与功耗。

上述前端加速工作表明，随着前端处理能力的大幅提升，后端优化环节的计算压力进一步加剧，因此针对后端的专用加速器设计变得尤为迫切。

\subsubsection{基于 FPGA 的后端加速器}

FPGA 凭借其高度可重构性、确定性时延以及相比 ASIC 更短的开发周期，成为 SLAM 后端加速研究的主流平台。早期的探索性工作中，T{\"o}rtei Tertei 等人\cite{tertei2016ekf_fpga}于 2016 年在 Computers \& Electrical Engineering 期刊上报告了面向三维视觉 SLAM 的 EKF 块加速器 FPGA 设计，通过将 EKF 更新步骤中的矩阵运算模块化并映射到 FPGA 流水线，验证了滤波法后端在可编程逻辑上的实现可行性，为后续优化法后端的 FPGA 加速提供了早期参考。

在该方向上，最具影响力的早期工作是 Liu 等人的 $\pi$-BA 系列。Qin 等人\cite{qin2019piba}于 2019 年在 FCCM 会议上首次提出了嵌入式 FPGA SoC 上的 BA 加速引擎 $\pi$-BA，该工作针对 BA 中三维点共观测（Co-observation）分布的不均匀性进行了深入分析，设计了基于共观测优化的计算调度策略，使硬件能够跳过大量无效的零块运算。$\pi$-BA 在 Xilinx ZC706 FPGA 上采用 Levenberg-Marquardt（LM）优化方法进行验证，将 Schur 消元算法部署到可编程逻辑端进行硬件加速，线性方程组求解等步骤则交由处理器系统（PS）端完成，实现了软硬协同的 BA 加速。Liu 等人\cite{liu2020pi}随后在 IEEE Transactions on Computers 上发表了 $\pi$-BA 的期刊扩展版本，进一步深化了共观测分布理论，完善了完整的 PL 端 BA 部署方案——包括利用非线性最小二乘（NLS）求解器计算 Jacobian 矩阵、通过 Schur 消元简化线性系统、以及通过 Cholesky 分解进行最终求解。通过硬件模块实现流水线并行与动态优化控制（动态调节 NLS 迭代次数和 Schur 消元模块数量），该工作在高效利用计算资源与提升系统效率方面取得了显著进展。

上海交通大学的 Sun 等人\cite{sun2020bax}在 2020 年提出了 BA 加速器 BAX，创新性地采用了访存/执行分离（Decoupled Access/Execute）的体系结构来缓解 BA 计算中严重的访存瓶颈。BAX 将内存访问与计算流水线分离到独立数据通路中，通过片内存储访问单元（Memory Access Unit, MAU）预取数据并分层缓存至计算流水线，同时设计了 3 种无转置矩阵乘法的数据流模式以减少矩阵转置操作带来的额外访存开销。在 Xilinx ZCU102 FPGA 上，BAX 在 200~MHz 频率下功耗仅为 1.12~W，相比 Intel i7-8700 桌面 CPU 实现了 1.73 倍的速度提升，相比 ARM Cortex-A53 嵌入式 CPU 则获得了 22.38 倍的加速。

麻省理工学院（MIT）的 Vivienne Sze 团队在低功耗视觉导航芯片设计方面做出了开创性工作。其代表作 Navion~\cite{suleiman2019navion} 是一款集成视觉惯性里程计（VIO）加速的 65~nm 专用 ASIC 芯片，功耗仅 24~mW。该芯片针对非线性优化中的稀疏性进行了深度挖掘，设计了专用的 Gauss-Newton 求解器，并采用二维稀疏块存储结构，相比移动端 CPU 实现了 171 倍的能效提升。此前，Suleiman 等人\cite{suleiman2018navion}在 2018 年 VLSI Circuits Symposium 上报告了 Navion 的初始设计，展示了将完整 VIO 流水线（包括特征跟踪、IMU 积分、非线性优化与更新）集成到单芯片上的可行性。虽然 Navion 针对 ASIC 流片设计，但其稀疏性利用策略与存储优化思想对 FPGA 架构设计极具启发价值。

Gan 等人\cite{gan2021eudoxus}于 HPCA 2021 上发表的 Eudoxus 是对自动驾驶定位系统进行系统性计算特征分析的代表性工作。该工作对包括 ORB-SLAM2、VINS-Mono 在内的多种主流定位算法进行了深入的性能剖析（Profiling），量化揭示了后端优化中矩阵运算（特别是 Hessian 矩阵构建与 Schur 补求解）在总计算时间中占比高达 70\%--90\% 的事实，并据此提出了针对热点函数的专用加速器设计方法学。在此基础上，Liu 等人\cite{liu2021archytas}在 MICRO 2021 上提出了 Archytas 框架，将加速器综合（Accelerator Synthesis）与运行时动态优化相结合，能够根据不同 SLAM 算法的计算特征自动生成定制化的硬件加速器。Archytas 通过参数化的硬件模板库和编译器驱动的设计空间探索，显著降低了 SLAM 加速器的开发成本。Liu 等人\cite{liu2022cicc}在 CICC 2022 上进一步提出了运行时可重构的 FPGA 加速器，支持在不同定位算法之间动态切换，并在 IEEE Transactions on Computers\cite{liu2022reconfig} 上发表了期刊扩展版本。

加州大学伯克利分校（UC Berkeley）的 Sophia Shao 团队在 2025 年 ASPLOS 会议上提出的 SuperNoVA 架构~\cite{zhao2025supernova}代表了算法-硬件协同设计（Algorithm-Hardware Co-Design）的最新趋势。该工作针对 SLAM 后端的计算延迟瓶颈，提出了资源感知的稀疏矩阵计算数据流架构，专门优化了 BA 过程中的舒尔补操作。通过动态调度与流水线重叠技术，SuperNoVA 在 Xilinx ZCU102 平台上实现了 10~ms 级别的后端优化延迟，相比 NVIDIA Jetson TX2 嵌入式 GPU 降低了 3.2 倍延迟并节省了 68\% 的能耗。

伦敦帝国理工学院（Imperial College London）Christos-Savvas Bouganis 团队探索了基于高斯置信度传播（Gaussian Belief Propagation, GBP）的分布式优化架构。传统的 BA 求解器（如 LM 或 Gauss-Newton）采用集中式策略，需要构建并分解全局 Hessian 矩阵，其计算模式天然适合序列化处理但难以充分发挥 FPGA 的并行优势。Sharif 等人在 DATE 2024 会议上提出的 GBP 加速框架~\cite{sharif2024framework}，通过高度并行的消息传递机制替代传统的集中式求解器，将全局 BA 问题分解为局部子问题并行求解。在因子图的每个节点上，GBP 通过迭代传递高斯分布形式的消息来达成全局一致的状态估计。该框架在 Xilinx Zynq-7000 平台上实现了 5.8 倍加速比，展示了非传统优化算法在硬件加速领域的应用潜力。

多伦多大学的 Chen 和 Anderson~\cite{chen2021fpga}则专注于稀疏 Cholesky 分解的硬件加速，这是 BA 求解器中的核心子模块。SLAM 后端优化中的正规方程 $\mathbf{H}\Delta\mathbf{x} = \mathbf{b}$ 通常通过 Cholesky 分解 $\mathbf{H} = \mathbf{L}\mathbf{L}^\mathrm{T}$ 求解。针对 SLAM 优化过程中稀疏模式动态变化的特点（不同时刻的共视关系导致 Hessian 矩阵的稀疏结构不断改变），Chen 等人提出了基于运行时任务调度的自适应求解器架构，通过符号分解预处理与数值分解流水化，在 Intel Arria~10 FPGA 上实现了相比 Intel MKL 库 2.1 倍的加速。该工作为处理不规则稀疏矩阵运算提供了硬件调度层面的解决方案。

此外，香港城市大学 Cheung 团队的郭润洲等人~\cite{guo2022fpga}针对双目视觉里程计的局部 BA 优化，在 Xilinx Zynq UltraScale+ MPSoC 上设计了专用加速器。该工作将矩阵乘加速器与预条件共轭梯度（PCG）求解器相结合，实现了 Jacobian 矩阵生成与 Hessian 矩阵累积的流水线并行，相比 Intel x86 处理器实现了约 1.74 倍的速度提升，相比 ARM 处理器获得了 27.76 倍的加速比，功耗节省达 95\%，在 MH\_03\_medium 数据集上达到了约 13.6~fps 的后端优化频率。

\subsubsection{基于 ASIC 的专用芯片}

与 FPGA 相比，ASIC 通过全定制化设计能够实现更高的能效比与更小的芯片面积，适合面向量产的嵌入式应用场景。除前述 MIT 的 Navion\cite{suleiman2019navion} 之外，Yoon 等人\cite{yoon2013graphics}于 2013 年在 IEEE TVLSI 上报告了一款面向 AR 应用的视觉与图形统一处理器，采用 0.18~$\mu$m CMOS 工艺，集成了姿态估计引擎，在 200~MHz 工作频率下达到 153.6~GOPS 的峰值算力，功耗仅为 413~mW，能效达 371.9~GOPS/W。该工作采用正交迭代（Orthogonal Iteration, OI）算法和四元数矩阵（Quaternion Matrices, QMs）来估计物体的旋转和平移，是较早将视觉位姿估计集成到专用芯片的尝试。

\subsubsection{基于 GPU 的并行优化}

GPU 凭借其大规模 SIMD 并行架构，在 BA 求解的矩阵运算中具有天然优势。Gopinath 等人\cite{gopinath2023gpu}在 ICRA 2023 上报告了针对视觉惯性 SLAM 局部 BA 的 GPU 加速工作，通过精细的 CUDA 内核设计与显存布局优化，在 NVIDIA 嵌入式 GPU 上实现了相比 CPU 实现数倍的加速。然而，GPU 方案的功耗通常在数瓦到数十瓦量级，难以满足无人机、可穿戴 AR 设备等功耗预算在百毫瓦级别的严苛场景，这也是 FPGA 与 ASIC 方案持续受到关注的根本原因。

\subsubsection{全系统软硬协同设计}

近年来，学术界开始关注覆盖 SLAM 前后端的全系统级硬件加速。Arora 等人\cite{arora2025socslam}在 DSD 2025 上提出的 SoC-SLAM 是这一方向的代表性工作，该系统在 Zynq FPGA SoC 上同时加速了前端特征提取/匹配与后端 BA 优化，通过 ARM 处理器与 FPGA 可编程逻辑的紧密协同，实现了完整视觉 SLAM 流水线的实时运行。Huang 等人\cite{huang2023eds}在 ICCAD 2023 上提出的 EDS-SLAM 则面向稠密双目 SLAM 场景，将基于深度学习的特征匹配（Learned Feature Matching）引入硬件加速架构，在保持高精度的同时实现了低功耗运行，代表了传统几何方法与深度学习方法在硬件层面融合的新趋势。

\subsection{国内研究现状}

近年来，国内高校在 SLAM 算法硬件化与后端加速器设计方面取得了长足进步，特别是在软硬协同设计、稀疏矩阵优化与能效提升方面涌现出多项高水平成果。与此同时，前端视觉处理模块的 FPGA 加速也受到了越来越多的关注。

在 SLAM 前端硬件加速方面，电子科技大学周军团队的 Liu 等人\cite{liu2022mobilesp}于 2022 年在 IEEE TCAS-I 上发表了 MobileSP 加速器，这是一款面向移动视觉 SLAM 的实时关键点提取硬件加速器。MobileSP 针对嵌入式平台的资源约束，设计了轻量级的可学习关键点检测网络及其对应的 FPGA 硬件映射方案，通过网络结构搜索与定点量化优化，在 FPGA 上实现了实时的关键点检测与描述子生成，为 SLAM 前端从传统手工特征（ORB/FAST）向深度学习特征的嵌入式部署迈出了重要一步。同一团队的 Liu 等人\cite{liu2024optflow_fpga}在 ISCAS 2024 上进一步提出了面向自动驾驶的超高性能可扩展光流 FPGA 加速器，通过多通道并行光流估计与可扩展的硬件架构设计，实现了满足高帧率视频处理需求的实时推理，展现了电子科技大学团队从后端优化向前端视觉处理全面拓展的技术布局。清华大学集成电路学院的 Zhang 等人\cite{zhang2022ropyslam}于 ICTA 2022 上提出了 ROPY-SLAM 异构 CPU-FPGA 系统，将前端特征提取与后端位姿优化分别部署在 FPGA 可编程逻辑与 ARM 处理器上，实现了软硬协同的完整 SLAM 流水线，是国内较早的覆盖 SLAM 前后端的全系统级硬件加速工作。

在 BA 优化硬件加速方向，清华大学集成电路学院张春团队保持了持续的系统性创新。年成等人\cite{nian2023energy}于 MWSCAS 2023 上首先提出了一种面向移动视觉 SLAM 机器人的节能位姿估计 FPGA 加速器，创新性地设计了基于 Cramer 定理和 Laplace 展开的六维线性方程求解器，在 Xilinx Kintex XCKU040 FPGA 上相比 Intel i7-10875H 处理器实现了 14.55 倍的速度提升和 45.63 倍的能效改善。同期，Niu 等人\cite{niu2023hwsw}在 IECON 2023 会议上提出了软硬协同矩阵求解方案，深入分析了非线性优化中 Levenberg-Marquardt 算法的计算特性，提出了 FSFI-Cholesky（基于 Cholesky 分解的直接解法，适用于稠密矩阵）和 FI-Iterative（基于迭代的求解方法，适合大规模稀疏矩阵）两种矩阵求解方法，在 ZCU104 FPGA 上分别实现了 12.8 倍与 120.2 倍的加速提升。在此基础上，年成等人\cite{nian2025vslamba}在 IEEE TVLSI 2025 上发表了 VSLAM-BA 加速器的期刊完整版本，提出了单帧束调整（Single-Frame Bundle Adjustment）架构，通过细粒度的计算重组与混合精度定点运算策略，实现了每帧仅 197~$\mu$J 的超低能耗，将乘法器资源使用量降低了 42\%，为移动机器人平台提供了极具竞争力的能效解决方案。清华大学电子工程系的另一研究团队也在相关方向取得了进展，Xiang 等人\cite{xiang2024loadbalanced}在 ICFPT 2024 上提出了面向基于 CNN 的视觉惯性里程计的负载均衡稀疏 BA 加速器，针对深度学习前端输出的特征点分布不均匀问题，设计了动态负载均衡的硬件调度机制，显著提升了各计算单元的利用率。莫霄睿等人\cite{mo2024ba_review}在 2024 年发表的综述论文《视觉 SLAM 机器人中光束法平差优化芯片研究综述》系统梳理了 BA 优化专用芯片的研究现状与发展趋势，涵盖了 FPGA、ASIC 与 GPU 三大硬件平台上的加速方法，为该领域的后续研究提供了重要参考。

在后端优化的另一条技术路线上，电子科技大学信息与通信工程学院周军团队形成了独特的研究体系。Liu 等人\cite{liu2023boha}于 2023 年在 IEEE TCAS-II 上发表了 BOHA（Backend Optimization Hardware Accelerator），提出了递归细粒度 H 矩阵分解（Recursive Fine-Grain H-Matrix Decomposition）策略，通过避免后端优化中的数据冗余来减小计算延迟，同时采用两级深度舒尔消元法，在 Xilinx FPGA ZCU104 上实现了 5.1~ms 的后端优化计算延迟和 0.27\% 的轨迹误差。在此基础上，黄坤\cite{huangkun2025vslam}在其硕士学位论文中系统总结了面向视觉 SLAM 的低延时位姿优化硬件加速器设计方法，提出了硬件友好型位姿优化算法与两阶段求解的多粒度流水架构，在 ZCU104 FPGA 上以 202.82~帧/秒的处理速率运行，轨迹误差仅为 0.26\%（无回环）和 0.15\%（有回环），展现了优异的延时-精度权衡。Liu 等人\cite{liu2025vslamba}于 2025 年在 IEEE TCAS-I 上进一步发表了 VSLAM-BA 系统，提出了算法与硬件协同设计的方法论，从算法层面重新审视了 BA 计算的数据依赖关系，据此设计了高性能且节能的后端硬件加速架构，代表了该团队在该方向的最新进展。

在 VIO 算法系统与硬件加速的结合方面，香港中文大学沈劭劼团队提出的 VINS-Mono 算法~\cite{qin2018vins}已成为学术界与工业界广泛采用的开源基准，为后续大量的算法改进与硬件加速工作提供了标准化的算法框架与评测基准。基于该算法框架，Li 等人~\cite{li2022fpga}在 IROS 2022 会议上针对后端优化瓶颈设计了 FPGA 加速器，采用双级并行架构加速 Jacobian 矩阵与残差向量的计算，相比 ARM Cortex-A53 处理器获得了 23 倍的处理频率提升。

综上所述，国内外研究均表明，随着 SLAM 系统对实时性与精度要求的不断提高，单纯依靠通用处理器的软件优化已难以满足需求。通过 FPGA 等专用硬件对后端非线性优化、稀疏矩阵构建与求解进行深度加速，已成为当前学术界与工业界共同关注的热点方向。表~\ref{tab:hw_accel_comparison} 对上述代表性工作进行了系统对比。

\begin{table}[htbp]
\centering
\bicaption{SLAM 后端硬件加速器代表性工作对比}{Comparison of representative SLAM backend hardware accelerators}
\label{tab:hw_accel_comparison}
\resizebox{\linewidth}{!}{
\begin{tabular}{lcccccc}
\hline
\textbf{工作} & \textbf{年份} & \textbf{平台} & \textbf{目标算法} & \textbf{求解方法} & \textbf{加速比} & \textbf{功耗} \\
\hline
$\pi$-BA\cite{liu2020pi} & 2020 & Xilinx ZC706 & BA (LM) & Schur + Cholesky & -- & -- \\
BAX\cite{sun2020bax} & 2020 & Xilinx ZCU102 & BA & Schur + 直接法 & 22.38$\times$ (vs ARM) & 1.12 W \\
Navion\cite{suleiman2019navion} & 2019 & 65 nm ASIC & VIO & Gauss-Newton & 171$\times$ (能效) & 24 mW \\
Guo et al.\cite{guo2022fpga} & 2022 & Zynq US+ & LBA & PCG & 27.76$\times$ (vs ARM) & -- \\
BOHA\cite{liu2023boha} & 2023 & ZCU104 & BA & Schur (两级) & -- & -- \\
Nian et al.\cite{nian2025vslamba} & 2025 & FPGA & SF-BA & Cramer/Laplace & 14.55$\times$ (vs i7) & 197 $\mu$J/帧 \\
SuperNoVA\cite{zhao2025supernova} & 2025 & ZCU102 & BA & Schur + 动态调度 & 3.2$\times$ (vs Jetson) & -- \\
GBP\cite{sharif2024framework} & 2024 & Zynq-7000 & BA (GBP) & 消息传递 & 5.8$\times$ & -- \\
\textbf{本文} & 2026 & Zynq US+ & SchurVINS & Schur + EKF & 3.74$\times$ (vs ARM) & -- \\
\hline
\end{tabular}
}
\end{table}

从表~\ref{tab:hw_accel_comparison} 可以看出，现有工作虽然在各自的技术路线上取得了显著进展，但普遍存在以下不足：（1）多数工作仅针对 BA 计算链中的单一阶段（如仅加速 Hessian 构建或仅加速线性求解），缺乏从 Jacobian 计算到 Schur 消元的全链路流水优化；（2）对 SLAM 后端特有的稀疏观测模式利用不够充分，未能充分跳过无效的零块计算；（3）存储优化策略较为粗糙，未能针对 Hessian 矩阵的对称性与块稀疏性设计紧凑的片上存储方案。本文针对上述问题，提出了基于 SOB 编码的动态门控机制与紧凑型 BRAM 管理方案，旨在实现更高效的全链路后端加速。

\section{本文主要研究内容与贡献}

针对上述国内外研究中存在的不足——全链路流水优化缺失、稀疏观测模式利用不充分以及片上存储方案粗糙等问题，本文围绕嵌入式平台下视觉 SLAM 后端优化计算负担重、实时性差这一核心矛盾，提出并系统实现了一套基于 FPGA 的软硬协同 SLAM 后端加速方案。本文以 SchurVINS 算法框架为基础，将"舒尔消元 + EKF 增量更新"的轻量化求解策略映射为可在 Zynq UltraScale+ 平台高效运行的异构硬件架构，从算法理论、系统架构、硬件模块设计到全系统集成验证形成了完整的研究体系。主要研究内容与贡献概括如下：

\begin{enumerate}
    \item \textbf{提出了面向嵌入式实时约束的"舒尔消元 + EKF 增量更新"算法策略}。针对传统 BA 在嵌入式平台上 LM 多轮迭代计算量过大的问题，本文将舒尔补消元后的等效信息矩阵与残差向量代入 EKF 观测更新步骤，以单次卡尔曼增益计算替代多轮矩阵分解迭代，将后端计算主体归结为高度规则化的雅可比累加与 $3 \times 3$ 矩阵求逆两个计算核，实现了算法复杂度与硬件友好性的统一，为后续硬件加速器的架构映射奠定了算法基础。

    \item \textbf{设计了基于 SOB 编码的两级动态门控机制与紧凑型存储管理方案}。针对 SLAM 后端 Hessian 矩阵构建中大量零块导致的无效计算问题，设计了一种 8 位的稀疏观测位图（Sparse Observation Bitmap, SOB）编码机制，紧凑表达滑动窗口内 4 个状态与左右目双相机的观测活跃拓扑。基于该编码，在硬件流水线前端实现了状态级旁路（State-level Bypass）与相机级时钟门控（Camera-level Clock Gating）两级稀疏跳过机制，自动跳过零块计算。SOB 解码得到的非零子块紧凑偏移直接用于 BRAM 写地址生成，实现了仅存储非零 Hessian 块的紧凑存储结构与隐式索引寻址。实验表明，该机制在典型稀疏率 50\% 场景下将总浮点运算量降低约 45\%，将 Hessian 矩阵的片上 BRAM 占用压缩至仅 67~KB，为同类方案中最优。

    \item \textbf{构建了支持双目复用的 5 级深度流水线雅可比计算模块与并行化 Schur 消元核心}。雅可比计算模块采用从世界坐标变换到 Hessian 矩阵在线累加的完整 5 级流水线架构，在坐标变换级通过切换外参常量实现左右目分支的单硬件路径复用，有效减少了资源占用；在残差级集成 Huber 鲁棒加权与观测方差归一化；在累加级利用 SOB 偏移实现非零子块的稀疏直写。Schur 消元模块面向路标块均为 $3 \times 3$ 对称正定矩阵的特点，以代数余子式法并行计算逆矩阵，引入奇异性保护机制防止数值溢出，并利用结果矩阵的对称性将输出写操作减少约 50\%。

    \item \textbf{完成了全系统集成、多数据集实验验证与真实场景部署测试}。基于 Zynq UltraScale+ 平台搭建了完整的软硬协同加速系统，设计了三通道 AXI 通信架构（AXI4-Lite 控制平面、AXI4-MM 数据平面与硬件中断通知）与双缓冲乒乓调度策略，有效隐藏了片外存储访问延迟。在 EuRoC 和 TUM-VI 权威数据集上的系统性实验表明，所设计的 Schur Kernel 将单次后端优化延迟从 2.84~ms 压缩至 0.76~ms，实现了 3.74 倍的加速比，系统后端优化频率稳定在 25~Hz 以上；EuRoC 上的平均绝对轨迹误差（ATE）为 0.073~m，与全软件基准（0.075~m）精度等价；单帧能耗仅为 19.1~mJ，为所有参照方案中最低。户外真实场景的部署测试进一步验证了系统在动态环境下的工程可行性。
\end{enumerate}

\section{论文结构安排}

本文共分六章，各章内容的逻辑关系与章节安排如图~\ref{fig:thesis_structure} 所示。

\begin{figure}[h!t]
\centering
\resizebox{0.98\linewidth}{!}{% thesis_structure.tex —— 论文结构安排横向树形图
\begin{tikzpicture}[
    >=Stealth,
    every node/.style={font=\rmfamily\footnotesize},
    % 根节点
    root/.style={
        draw, rounded corners=4pt, thick, align=center,
        fill=blue!18, minimum width=24mm, minimum height=14mm,
        text width=22mm
    },
    % 章节节点
    chap/.style={
        draw, rounded corners=3pt, thick, align=center,
        fill=blue!8, minimum width=34mm, minimum height=8mm,
        text width=32mm, font=\rmfamily\footnotesize
    },
    % 内容叶节点
    leaf/.style={
        draw, rounded corners=2pt, thin, align=left,
        fill=gray!6, minimum width=42mm, minimum height=6mm,
        text width=40mm, font=\rmfamily\scriptsize
    },
    % 连线
    edge/.style={-, thick, gray!70},
    branch/.style={-, thin, gray!50},
]

%% ── 根节点 ───────────────────────────────────────────
\node[root] (root) at (0, 0)
    {\textbf{面向嵌入式异构}\\[-1pt]\textbf{系统的 SLAM 算法}\\[-1pt]\textbf{软硬件协同加速设计}};

%% ── 章节节点（纵向均布，列 x=4.0cm）───────────────────
\def\cs{2.2}  % 章间距

\node[chap] (c1) at (4.0,  5*\cs/2) {\textbf{第一章}\ \ 绪\ \ 论};
\node[chap] (c2) at (4.0,  3*\cs/2) {\textbf{第二章}\ \ 算法理论基础};
\node[chap] (c3) at (4.0,  1*\cs/2) {\textbf{第三章}\ \ 系统架构总体设计};
\node[chap] (c4) at (4.0, -1*\cs/2) {\textbf{第四章}\ \ 硬件加速器详细设计};
\node[chap] (c5) at (4.0, -3*\cs/2) {\textbf{第五章}\ \ 系统实现与实验评价};
\node[chap] (c6) at (4.0, -5*\cs/2) {\textbf{第六章}\ \ 总结与展望};

%% ── 根 → 章节连线（主干 + 纵向脊 + 水平分支）──────────
\def\rx{2.0}  % 根节点主干分叉脊的 x 坐标

\draw[edge] (root.east) -- (\rx, 0);
\draw[edge] (\rx, 5.5) -- (\rx, -5.5);
\foreach \c/\cy in {c1/5.5, c2/3.3, c3/1.1, c4/-1.1, c5/-3.3, c6/-5.5}{
    \draw[edge] (\rx, \cy) -- (\c.west);
}

%% ── 内容叶节点（列 x=9.2cm） ──────────────────────────
\def\lx{9.2}

% 第一章
\node[leaf] (l1a) at (\lx,  5*\cs/2 + 0.65) {研究背景与应用意义};
\node[leaf] (l1b) at (\lx,  5*\cs/2 - 0.00) {国内外研究现状综述};
\node[leaf] (l1c) at (\lx,  5*\cs/2 - 0.65) {本文主要贡献与论文结构};

% 第二章
\node[leaf] (l2a) at (\lx,  3*\cs/2 + 0.65) {相机模型与视觉前端处理};
\node[leaf] (l2b) at (\lx,  3*\cs/2 - 0.00) {IMU 预积分与紧耦合状态估计};
\node[leaf] (l2c) at (\lx,  3*\cs/2 - 0.65) {BA 稀疏性、舒尔补与 EKF 更新};

% 第三章
\node[leaf] (l3a) at (\lx,  1*\cs/2 + 0.50) {设计目标与软硬件功能划分};
\node[leaf] (l3b) at (\lx,  1*\cs/2 - 0.50) {AXI 通信协议与双缓冲调度};

% 第四章
\node[leaf] (l4a) at (\lx, -1*\cs/2 + 0.65) {五级深度流水线雅可比计算模块};
\node[leaf] (l4b) at (\lx, -1*\cs/2 - 0.00) {SOB 编码与两级动态门控机制};
\node[leaf] (l4c) at (\lx, -1*\cs/2 - 0.65) {3$\times$3 并行 Schur 消元核心};

% 第五章
\node[leaf] (l5a) at (\lx, -3*\cs/2 + 0.65) {传感器标定与实验环境搭建};
\node[leaf] (l5b) at (\lx, -3*\cs/2 - 0.00) {轨迹精度与计算延迟分析};
\node[leaf] (l5c) at (\lx, -3*\cs/2 - 0.65) {资源利用率、能效与 SOB 优化};

% 第六章
\node[leaf] (l6a) at (\lx, -5*\cs/2 + 0.40) {全文总结与主要创新点};
\node[leaf] (l6b) at (\lx, -5*\cs/2 - 0.40) {后续工作展望};

%% ── 章节 → 叶节点连线（主干 + 垂直脊 + 水平分支）────────
\def\mx{6.8}  % 垂直分支脊的 x 坐标

% 第一章：y=5.5，叶偏移 ±0.65、0
\draw[branch] (c1.east) -- (\mx, 5.5);
\draw[branch] (\mx, 5.5+0.65) -- (\mx, 5.5-0.65);
\draw[branch] (\mx, 5.5+0.65) -- (l1a.west);
\draw[branch] (\mx, 5.5+0.00) -- (l1b.west);
\draw[branch] (\mx, 5.5-0.65) -- (l1c.west);

% 第二章：y=3.3，叶偏移 ±0.65、0
\draw[branch] (c2.east) -- (\mx, 3.3);
\draw[branch] (\mx, 3.3+0.65) -- (\mx, 3.3-0.65);
\draw[branch] (\mx, 3.3+0.65) -- (l2a.west);
\draw[branch] (\mx, 3.3+0.00) -- (l2b.west);
\draw[branch] (\mx, 3.3-0.65) -- (l2c.west);

% 第三章：y=1.1，叶偏移 ±0.50
\draw[branch] (c3.east) -- (\mx, 1.1);
\draw[branch] (\mx, 1.1+0.50) -- (\mx, 1.1-0.50);
\draw[branch] (\mx, 1.1+0.50) -- (l3a.west);
\draw[branch] (\mx, 1.1-0.50) -- (l3b.west);

% 第四章：y=-1.1，叶偏移 ±0.65、0
\draw[branch] (c4.east) -- (\mx, -1.1);
\draw[branch] (\mx, -1.1+0.65) -- (\mx, -1.1-0.65);
\draw[branch] (\mx, -1.1+0.65) -- (l4a.west);
\draw[branch] (\mx, -1.1+0.00) -- (l4b.west);
\draw[branch] (\mx, -1.1-0.65) -- (l4c.west);

% 第五章：y=-3.3，叶偏移 ±0.65、0
\draw[branch] (c5.east) -- (\mx, -3.3);
\draw[branch] (\mx, -3.3+0.65) -- (\mx, -3.3-0.65);
\draw[branch] (\mx, -3.3+0.65) -- (l5a.west);
\draw[branch] (\mx, -3.3+0.00) -- (l5b.west);
\draw[branch] (\mx, -3.3-0.65) -- (l5c.west);

% 第六章：y=-5.5，叶偏移 ±0.40
\draw[branch] (c6.east) -- (\mx, -5.5);
\draw[branch] (\mx, -5.5+0.40) -- (\mx, -5.5-0.40);
\draw[branch] (\mx, -5.5+0.40) -- (l6a.west);
\draw[branch] (\mx, -5.5-0.40) -- (l6b.west);

\end{tikzpicture}
}
\bicaption{论文结构安排示意图}{Illustration of thesis structure}
\label{fig:thesis_structure}
\end{figure}

\textbf{第一章}为绪论，介绍研究背景与工程意义，梳理 SLAM 前后端算法的国内外研究现状及现有硬件加速工作的局限性，提出本文的研究目标，并概括全文主要贡献与论文结构安排。

\textbf{第二章}为算法理论基础，系统建立本文软硬协同后端所依赖的数学理论体系。以针孔相机模型与畸变校正为起点，介绍基于 SVO 2.0 半直接法的视觉前端处理流程；推导 IMU 预积分理论与紧耦合视觉惯性状态估计框架；深入分析 Hessian 矩阵的块稀疏结构，推导舒尔补消元降维公式，并提出面向嵌入式实时约束的"舒尔消元 + EKF 增量更新"轻量化求解策略。

\textbf{第三章}为系统架构总体设计，从顶层确立系统的四项量化设计目标，遵循"逻辑密集型任务归 PS、数据密集型规则计算归 PL"的原则完成软硬件功能划分，并设计三通道 AXI 通信协议（AXI-Lite 控制平面、AXI-MM 数据平面、硬件中断通知）与双缓冲乒乓调度策略，为第四章硬件模块设计提供系统级框架。

\textbf{第四章}为硬件加速器详细设计，是全文核心章节。依次介绍支持左右目复用的五级深度流水线雅可比计算模块、基于 8 位 SOB 编码的状态级旁路与相机级时钟门控两级动态门控机制、仅存储非零 Hessian 子块的紧凑 BRAM 管理方案，以及代数余子式法并行化 $3\times3$ 对称矩阵求逆的 Schur 消元核心。

\textbf{第五章}为系统实现与实验评价，详细记录基于 Zynq UltraScale+ 平台的全系统集成过程，包括三阶段传感器离线标定方法与结果。通过在 EuRoC 和 TUM-VI 权威数据集及真实场景上的系统性实验，从轨迹精度（ATE）、计算延迟与加速比、硬件资源占用、功耗与能效，以及 SOB 编码优化效果五个维度对加速器性能进行全面定量评估。

\textbf{第六章}为总结与展望，对全文研究工作进行系统回顾，提炼四项主要创新点，并就前端硬件化集成、广角相机模型扩展、回环检测全局优化加速及语义 SLAM 软硬协同加速等方向给出后续工作展望。

